\begin{theorem}[零点块对持久大纤维偏置的无力性]\label{thm:xi-time-part9zoa-zero-block-maxfiber-impotence}
\leanverified{paper\_xi\_time\_part9zoa\_zero\_block\_maxfiber\_impotence}
沿用定理 \ref{thm:xi-time-part9zo-zero-block-macroscopic-normalized-silent}、
定理 \ref{thm:xi-time-part60ab-window6-zero-block-halfscale-information-improvement}
与定理 \ref{thm:xi-time-part9m-maxfiber-uniformity-gap-chain} 的记号。若
\[
M:=F_{m+2},
\qquad
\mathcal Z_m:=\{k\in \ZZ/(M\ZZ):\ \widehat d_m(k)=0\},
\qquad
\eta_m:=\frac{|\mathcal Z_m|}{M},
\]
并且
\[
m+2\equiv 0\pmod 6,
\]
则虽有确定性的半阶零点块下界
\[
|\mathcal Z_m|\ge F_{(m+2)/2},
\]
但仍满足
\[
\eta_m=O(\varphi^{-m/2}).
\]
更具体地，任何仅通过
\[
M-|\mathcal Z_m|,
\qquad
\log(M-|\mathcal Z_m|),
\qquad
\sqrt{M-1-|\mathcal Z_m|}
\]
进入的零点驱动上界，其相对于零点自由基线的改变量都至多具有 \(O(\eta_m)=O(\varphi^{-m/2})\) 的相对尺度。与此相对，最大纤维偏置满足
\[
\liminf_{m\to\infty}\frac{M_m}{\overline d_m}\ge 1+C_\varphi>1.
\]
因此，算术零点块无法消除持久的 \(L^\infty\) 级非均匀性；真正顽固的障碍并不在零点集合，而在大纤维偏置本身。
\end{theorem}
\begin{proof}
由定理 \ref{thm:xi-time-part9zo-zero-block-macroscopic-normalized-silent}，
\[
\eta_m=\frac{|\mathcal Z_m|}{M}
=
\tau(m+2)\,O(\varphi^{-m/2})
=
O(\varphi^{-m/2}),
\]
其中最后一步仅用到约数函数的次幂增长远慢于任意指数函数。

同一定理还给出
\[
M-|\mathcal Z_m|=M(1-\eta_m),
\]
\[
\log(M-|\mathcal Z_m|)=\log M+O(\eta_m),
\]
\[
\sqrt{M-1-|\mathcal Z_m|}
=
\sqrt M\bigl(1+O(\eta_m)\bigr).
\]
故凡是仅通过这三类表达式进入的 zero-driven 上界，其与零点自由主项相比都只能产生 \(O(\eta_m)\) 级的相对改动。这正是“零点块只能提供半指数级副修正”的精确形式。

另一方面，定理 \ref{thm:xi-time-part9m-maxfiber-uniformity-gap-chain} 已经给出
\[
\liminf_{m\to\infty}\frac{M_m}{\overline d_m}\ge 1+C_\varphi>1.
\]
因此最大纤维相对均值的偏置在极限中仍与 \(1\) 保持严格正距离，不可能被任何 \(O(\varphi^{-m/2})\) 级的 zero-block 修正所抹平。于是折叠分布的持久非均匀性并非由零点块驱动，真正的主障碍只能来自大纤维偏置。证毕。
\end{proof}

\endinput
